\documentclass[10pt,a4paper]{amsart}
\usepackage[margin=19mm]{geometry}
\usepackage{amsmath,amssymb,mathtools}
\usepackage[scr=rsfs,cal=euler]{mathalfa}
\usepackage{xfrac}
\usepackage[unicode,hidelinks,bookmarks=false]{hyperref}
\newcommand{\ie}{i.e.\ }
\newcommand{\cf}{cf.\ }
\newcommand{\resp}{respectively\ }
\newcommand{\Subsection}{\subsection}
\providecommand{\nobreakdash}{\nobreak\hbox{-}}
\setlength{\emergencystretch}{2em}
\multlinegap0pt
\usepackage{mathtools}
\usepackage[shortlabels]{enumitem}
\newtheorem{teo}{Theorem}
\newtheorem{lem}{Lemma}
\newcommand{\N}{\mathbb{N}}
\newcommand{\R}{\mathbb{R}}
\newcommand{\bv}{\operatorname{BV}}
\newcommand{\leb}{\mathcal{L}}
\newcommand{\nl}{\left\|}
\newcommand{\nr}{\right\|}
\newcommand{\p}{{\mathfrak{p}}}
\newcommand{\pp}{\partial}
\newcommand{\res} {\mathop{\hbox{\vrule height 7pt width .5pt depth 0pt \vrule height .5pt width 6pt depth 0pt}}\nolimits}
\newcommand{\s}{\mathcal{S}}
\newcommand{\se}{\subseteq}
\newcommand{\vp}{\varphi}
\title{Intrinsic perimeter, compactness and Poincaré inequality for SBV functions in Carnot-Carathéodory spaces}
\author{Marco Di Marco}
\date{Source: arXiv:2510.19380v2 (CC BY 4.0)}
\begin{document}
\maketitle

\noindent\emph{Source and licence.} Intrinsic perimeter, compactness and Poincaré inequality for SBV functions in Carnot-Carathéodory spaces. Version arXiv:2510.19380v2 (CC BY 4.0), distributed by arXiv under the Creative Commons Attribution 4.0 International licence, http://creativecommons.org/licenses/by/4.0/, as recorded in the arXiv licence metadata for this exact version and shown on the abstract page https://arxiv.org/abs/2510.19380v2. Full licence notice: \texttt{licenses/arxiv-2510.19380-CC-BY-4.0.txt}.\par\smallskip

\noindent\textit{Editorial scope.} Counted unit: the article's lem lemma\_chain (source0.tex lines 417--422). The article's complete original proof of it is delivered with it (source0.tex lines 423--454, one \texttt{proof} environment of 32 lines). No mathematics is written, repaired or re-derived: the statement and the proof are the article's own lines, in the article's own notation, and no retained line is substituted or re-flowed.  Retained uncounted as the source-local context the proof needs: lines 330--357.  Nothing was omitted from any retained span, so the package carries no omission record.  The retained lines cite 2 works, whose entries are transcribed verbatim from the article's own bibliography source in the e-print and delivered with short editorial labels recorded in the item's reference mapping.  No retained line contains a figure environment, an image-inclusion command or drawing code, so no image is delivered and the package declares no asset dependency.  Editorial formatting only, in addition: an AMS \texttt{amsart} preamble that restates the article's own declarations and macros used by the retained lines, namely the environments \texttt{teo}, \texttt{lem} on separate counters, together with the standard packages those lines need.  The retained environments print in this document as Teo 1, Lemma 1 in order of appearance, whereas the article prints the same items with the numbering of its own full document.  The complete native source of the article as distributed in the arXiv e-print for this exact version is bundled unchanged as \texttt{sources/187566/source0.tex}.
\begin{teo}[{\cite[Theorem 2.3.5]{fssc}}]\label{teo_coarea1}
Let $u \in \bv_X(\Omega)$. Then $\{u>t\}$ has finite $X$-perimeter for $\leb^1$-a.e. $t \in \R$ and
\[
|D_Xu|(\Omega)=\int_{-\infty}^{+\infty} P_X(\{u >t\},\Omega)dt.
\]
Conversely, if $u \in L^1(\Omega)$ and 
\[
\int_{-\infty}^{+\infty} P_X(\{u >t\},\Omega)dt<+\infty,
\]
then $u \in \bv_X(\Omega)$.
\end{teo}
Let us observe that also the following version of the coarea formula holds. Although the proof is standard, we provide it for the sake of completeness.
\begin{teo}\label{teo_coarea2}
    Let $u \in \bv_X(\Omega)$. Then for every Borel set $A \se \Omega$ we have
\[
D_Xu(A)=\int_{-\infty}^{+\infty} D_X \chi_{\{u >t\}}(A) dt.
\]
\end{teo}
\begin{proof}
    Let $\vp \in C^1_c(\Omega)$ and $i \in \{1,\dots,m\}$. We have
    \begin{align*}
        \int_\Omega \vp dD_{X_i}u&=-\int_\Omega u(x)X^*_i\vp(x)dx\\
        &=-\int_\Omega \left( \int_0^{+\infty} \chi_{\{u>t\}}(x)dt \right)X^*_i\vp(x)dx+\int_\Omega \left( \int_{-\infty}^0 (1- \chi_{\{u>t\}})(x)dt \right)X^*_i\vp(x)dx \\
        &=-\int_0^{+\infty} \left( \int_\Omega \chi_{\{u>t\}}(x)X^*_i\vp(x)dx \right)dt+\int_{-\infty}^0 \left(\int_\Omega  (1- \chi_{\{u>t\}})(x)X^*_i\vp(x)dx \right)dt\\
        &= \int_{-\infty}^{+\infty} \left( \int_\Omega \vp dD_{X_i} \chi_{\{ u>t \} } \right)dt,
    \end{align*}
    where we used Fubini's theorem and the fact that, for every positive Borel function $v$, we have $v(x)=\int_0^{+\infty} \chi_{ \{ v>t \} }(x)dt$.
\end{proof}

\begin{lem}\label{lemma_chain}
    Let $u \in \bv_X(\Omega)$. Let $\phi \in C^1(\R) \cap W^{1,\infty}(\R)$. Then $ \phi \circ u\in \bv_X(\Omega)$ and
    \begin{equation}\label{eq_chainrule}
    D_X(\phi \circ u) \res (\Omega \setminus \s_u) =\phi'(u^\p) D_Xu \res (\Omega \setminus \s_u).
    \end{equation}
\end{lem}

\begin{proof}
    Let us first prove that $\phi \circ u \in \bv_X(\Omega)$. By \cite[Theorem 2.2.2]{fssc} there exists a sequence of functions $(u_k)_{k \in \N}$ such that, for every $k \in \N$, we have $u_k \in C^\infty(\Omega) \cap \bv_X(\Omega)$ and
    \[
    \nl u-u_k \nr_{L^1(\Omega)} \xrightarrow{k \to +\infty}0, \qquad |D_Xu_k|(\Omega) \xrightarrow{k \to +\infty}|D_Xu|(\Omega).
    \]
    Then $\nl \phi \circ u-\phi \circ u_k \nr_{L^1(\Omega)} \xrightarrow{k \to +\infty}0$ and, by the lower semicontinuity of the total variation, 
    \[
    |D_X(\phi \circ u)|(\Omega) \leq \liminf_{k \to +\infty} |D_X(\phi \circ u_k)|(\Omega) \leq \nl \phi' \nr_{\infty}\liminf_{k \to +\infty}|D_Xu_k|(\Omega)=\nl \phi' \nr_{\infty}|D_Xu|(\Omega)<\infty,
    \]
concluding $\phi \circ u \in \bv_X(\Omega)$. Since every $\phi \in C^1(\R) \cap W^{1,\infty}(\R)$ can be written as the difference of two strictly increasing functions with the same properties, in the rest of the proof we assume, without loss of generality\footnote{By the linearity of \eqref{eq_chainrule}.}, that $\phi$ is also strictly increasing and, in particular, injective. By the Lipschitzianity of $\phi$ we have $\s_{\phi \circ u} \se \s_u$. Let $A \se \Omega \setminus \s_u$ be any Borel set. By the coarea formula (Theorem \ref{teo_coarea2}) we have
\[
D_X(\phi \circ u)(A)=\int_{-\infty}^{+\infty} D_X\chi_{  \{ \phi \circ u >s  \}  } (A)ds.
\]
By the change of variable $s=\phi(t)$ we obtain
\[
\int_{-\infty}^{+\infty} D_X\chi_{\{\phi \circ u>s \}} (A)ds=\int_{-\infty}^{+\infty} \phi'(t)D_X\chi_{\{u>t \}} (A)dt.
\]
We have that if $x \in A \cap \pp^* \{u>t \} $, then $u^\p(x)=t$ (see for instance \cite[Proposition 2.22]{dv}) and 
\[
\int_{-\infty}^{+\infty} \phi'(t)D_X\chi_{\{u>t \}} (A)dt=\int_{-\infty}^{+\infty}\int_A\phi'(u^\p(x))dD_X\chi_{\{u>t \}}(x)dt
\]
Using again the coarea formula (Theorem \ref{teo_coarea2}) and Fubini's theorem we obtain
\[
\int_{-\infty}^{+\infty}\int_A\phi'(u^\p(x))dD_X\chi_{\{u>t \}}(x)dt=\int_A \phi'(u^\p)dD_Xu
\]
the latter implying that 
\[
D_X(\phi \circ u) \res (\Omega \setminus \s_u)=\phi'(u^\p)D_Xu\res (\Omega \setminus \s_u),   
\]
concluding the proof.

\end{proof}

\begin{thebibliography}{99}
\bibitem[dv]{dv} {\sc Don, S., and Vittone, D.} \newblock Fine properties of functions with bounded variation in {C}arnot-{C}arath\'{e}odory spaces. \newblock {\em J. Math. Anal. Appl. 479}, 1 (2019), 482--530.
\bibitem[fssc]{fssc} {\sc Franchi, B., Serapioni, R., and Serra~Cassano, F.} \newblock Meyers-{S}errin type theorems and relaxation of variational integrals depending on vector fields. \newblock {\em Houston J. Math. 22}, 4 (1996), 859--890.
\end{thebibliography}
\end{document}
